\documentclass[11pt]{letter}
\usepackage[T1]{fontenc}
\usepackage[utf8]{inputenc}
\usepackage{lmodern}
\usepackage[margin=1in]{geometry}
\usepackage{amsmath}
\usepackage{hyperref}

\signature{Tristan Simas\\McGill University}

\date{April 6, 2026}

\begin{document}

\begin{letter}{Editors\\
Journal of Computer and System Sciences}

\opening{Dear Editors,}

I submit the manuscript ``Zero-Error Recovery from Deterministic Partial Views'' for consideration in the Journal of Computer and System Sciences.

The paper develops a graph-theoretic theory of zero-error recovery from deterministic partial views of finite latent tuples. Admissible coordinate-subset views induce a confusability graph on latent states, exact recovery becomes equivalent to graph colorability, repeated composition lifts the one-shot graph to strong powers, and the normalized block rates converge to the Shannon capacity of the induced graph. On the full labeled tuple space, the realizable confusability relations are characterized exactly as those determined by upward-closed families of coordinate-agreement sets.

The manuscript also identifies a structural equality route for the induced graph class: when confusability is transitive, the graph collapses to a cluster graph, yielding capacity--$\vartheta$ equality, with meet-witnessing and fiber coherence as sufficient conditions. Under an affine restriction on the realized state family, the coordinate side carries a representable matroid whose rank gives tractable upper bounds on confusability and capacity. The main text focuses on the core graph-capacity theory, while the boundary proofs and representative host instantiations appear in appendices.

All theorem-level claims cited in the manuscript are machine-checked in Lean 4. An archived Zenodo record containing the formalization, manuscript PDF, and submission source bundle is available at \url{https://doi.org/10.5281/zenodo.19354075}. The manuscript is not under consideration elsewhere.

Thank you for your time and consideration.

\closing{Sincerely,}

\end{letter}
\end{document}
